%% The first command in your LaTeX source must be the \documentclass command.
%%
%% Options:
%% twocolumn : Two column layout. Do not use twocolumn for papers submitted to CEUR-WS!
%% hf: enable header and footer.
\documentclass[
% twocolumn,
% hf,
]{ceurart}

%%
%% One can fix some overfulls
\sloppy

%%
%% Minted listings support 
%% Need pygment <http://pygments.org/> <http://pypi.python.org/pypi/Pygments>
\usepackage{xcolor}
\usepackage{listings}

% --- JavaScript + Unravel-RDF DSL ------------------------------------
\lstdefinelanguage{JavaScript}{
  keywords={const, let, var, function, return, if, else, for, while,
            class, new, await, async},
  sensitive=true,
  comment=[l]{//},
  morecomment=[s]{/*}{*/},
  morestring=[b]",
  morestring=[b]'
}

\lstdefinestyle{unraveljs}{
  language=JavaScript,
  basicstyle=\ttfamily\small,
  % Mots-clés JS (const, etc.) en bleu
  keywordstyle=\color{blue}\bfseries,
  commentstyle=\color{gray}\itshape,
  stringstyle=\color{red!60!black},
  numbers=none,
  frame=single,
  breaklines=true,
  columns=fullflexible,
  % Mots-clés supplémentaires pour mettre en évidence l'API Unravel
  morekeywords={UnravelSession,something,from,out},
  % API Unravel en teal gras
  keywordstyle=[2]\color{teal}\bfseries,
  classoffset=1,
  otherkeywords={UnravelSession,something,from,out},
  classoffset=0,
  % Autoriser l'injection de LaTeX dans le code JS
  escapeinside={(*}{*)}
}

% --- SPARQL -----------------------------------------------------------
\lstdefinelanguage{SPARQL}{
  morekeywords={
    SELECT, CONSTRUCT, ASK, DESCRIBE,
    WHERE, PREFIX, BASE,
    FILTER, OPTIONAL, GRAPH, UNION,
    LIMIT, OFFSET, DISTINCT, REDUCED,
    ORDER, BY, GROUP, HAVING
  },
  sensitive=true,
  comment=[l]{\#},
  morestring=[b]"
}

\lstdefinestyle{sparqlstyle}{
  language=SPARQL,
  basicstyle=\ttfamily\small,
  keywordstyle=\color{blue}\bfseries,
  commentstyle=\color{gray}\itshape,
  stringstyle=\color{red!60!black},
  numbers=none,
  frame=single,
  breaklines=true,
  columns=fullflexible,
  escapeinside={(*}{*)}
}

% to active note:
%\usepackage{todonotes}
% to unactive note:
\usepackage[disable]{todonotes}

%%
%% end of the preamble, start of the body of the document source.
\begin{document}

%%
%% Rights management information.
%% CC-BY is default license.
\copyrightyear{2022}
\copyrightclause{Copyright for this paper by its authors.
  Use permitted under Creative Commons License Attribution 4.0
  International (CC BY 4.0).}

%%
%% This command is for the conference information
\conference{ISWC 2026}

%% Journal of Web Semantics
%%
%% We position Unravel-RDF not merely as a software artifact, but as a
%% developer-oriented abstraction layer intended to lower the barrier to
%% adopting Semantic Web technologies in interactive applications.

%% The "title" command
\title{Unravel-RDF: A Developer-Oriented Programming Model for Interactive
Applications over RDF Knowledge Graphs}

\tnotemark[1]

\author[1,4]{Olivier Filangi}[%
  orcid=0000-0002-2094-3271,
  email=olivier.filangi@inrae.fr,
]

\cormark[1]

\author[2,4]{Nils Paulhe}[%
  orcid=0000-0003-4550-1258
]

\author[3,4]{Guillaume Laisney}[%
  orcid=0009-0005-5653-3257
]

\author[3,4]{Clément Frainay}[%
  orcid=0000-0003-4313-2786
]

\author[2,4]{Franck Giacomoni}[%
  orcid=0000-0001-6063-4214
]

\address[1]{Institute for Genetics, Environment and Plant Protection (IGEPP),
INRAE, Université de Rennes, Institut Agro Rennes-Angers, Le Rheu, France}

\address[2]{Human Nutrition Unit (UNH), UMR 1019 INRAE–Université d’Auvergne,
Clermont-Ferrand, France}

\address[3]{Toxalim, Research Centre in Food Toxicology, INRAE,
Université de Toulouse, ENVT, INP-Purpan, Toulouse, France}

\address[4]{MetaboHUB National Infrastructure for Metabolomics and Fluxomics,
France}

%% Footnotes
\cortext[1]{Corresponding author.}


\begin{abstract}
  Hello World
\end{abstract}


\begin{keywords}
  RDF Knowledge graphs \sep
  Frontend development \sep
  Domain-specific language (DSL) \sep
  Metabolomics \sep
\end{keywords}

\maketitle

\section{Introduction}

%% comparaison avec les ORM relationnels ;
%%% pourquoi RDF n'a pas d'équivalent frontend ;
%%% pourquoi SPARQL brut ne scale pas dans une application.

Semantic Web technologies are widely used to represent and query knowledge in
the life sciences\citep{wu_semantic_2014}. Many resources are published as RDF
knowledge graphs and made available to the community via downloadable datasets
or SPARQL endpoints, thereby enabling interoperability between heterogeneous
datasets describing molecular entities, biological processes, pathways, and
phenotypic characteristics across diverse
organisms\citep{LODCloudLifeSciences}. Semantic Web technologies support the
development of bioinformatics tools through the integration and reuse of such
datasets, but their effective adoption within analysis platforms and web
applications still poses significant technical challenges.

Beyond SPARQL endpoints and RDF dumps, several tools provide visual or guided 
interfaces for exploring RDF knowledge graphs, including query-based faceted 
search (QFS) systems such as Sparklis\citep{ferre_sparklis_2016} and platforms 
like Sparnatural\citep{francart_sparnatural_2023}, Sampo-UI\citep{ikkala_sampo-ui_2021}, 
and SemTK\citep{cuddihy_semtk_2017}. In particular, QFS approaches introduce abstractions 
such as focus and progressive query refinement that enable users to navigate RDF 
data without writing SPARQL manually. However, these systems are primarily designed 
for end users through dedicated interfaces and do not expose their query abstractions 
as reusable, developer-oriented frameworks for building RDF applications.

In practice, frontend developers building web applications on top of RDF
knowledge graphs increasingly rely on the JavaScript ecosystem. The RDF
JavaScript Libraries Community Group -- Tools\footnote{\url{https://www.w3.org/groups/cg/rdfjs/tools/}}
has played a key role in this ecosystem by defining common interfaces and
representations that enable interoperability between RDF libraries and
engines. This ecosystem, including SPARQL engines such as
Comunica\citep{taelman_comunica_2018}, provides essential building blocks for
querying heterogeneous RDF sources directly from JavaScript applications.
However, these tools remain low-level infrastructure: they expose RDF
primitives and query execution mechanisms but do not provide
developer-oriented abstractions for structuring application data access, such
as composable query builders, domain-specific languages (DSLs), or
application-level data-access models commonly found in relational database
ecosystems.

As a result, there remains a significant gap between RDF query
formulation and modern frontend development practices, leading to
brittle integrations and weak alignment between knowledge graph queries
and application-level structures. To address this gap, we introduce
Unravel-RDF, a developer-oriented library for building applications on
top of RDF knowledge graphs. Rather than a user-facing visual interface
or a standalone query generator, Unravel-RDF combines a DSL for SPARQL
query construction with a composable API tailored to event-driven web
applications. Inspired by the notion of \emph{focus} introduced in
query-based faceted search (QFS) systems, Unravel-RDF elevates this
concept into a reusable programming abstraction for composing,
executing, and integrating SPARQL queries within application code. We
evaluate Unravel-RDF on the \textit{FORVM DiseaseChem}
\citep{delmas_forum_2021} and the experimental \textit{FORVM Plants}
knowledge graphs through an inter-domain exploration workflow deployed
as a web application, focusing on its usability for developers and its
support for structuring application-level access to RDF data.

\section{Background and Related Work}

\subsection{RDF application development}

%% Comparer :

%% Apache Jena (Java)
%% RDF4J (Java)
%% rdflib (Python)
%% RDFJS ecosystem (JavaScript)

%% Question :"What abstractions exist for developers?"

\subsection{RDF/JS and JavaScript Semantic Web ecosystem}

%% Détailler :

%% RDFJS interfaces ;
%% Comunica ;
%% Linked Data Fragments ;
%% Solid éventuellement.

%% Message: "These initiatives demonstrate the importance of Web-native access to
%% distributed RDF resources, but they do not address application-level query
%% composition."

\subsection{Visual query systems and QFS}

%% Là tu développes :

%% Sparklis ;
%% Sparnatural ;
%% Sampo ;
%% SemTK.

%% Puis :
%% 
%% We reuse the notion of focus, but move it from an end-user interaction
%% model to a developer programming abstraction.

\section{Architecture and Programming Model}



\subsection{Design principles}

%% expliques :

%% immutable objects ;
%% separation of concerns ;
%% composability ;
%% frontend integration.




\subsection{Layered Architecture}

Figure~\ref{fig:architecture} shows a three-layer architecture: an
\emph{Application Layer} for user-facing logic, an \emph{Unravel Query
Layer} providing the API for query construction, and an \emph{Execution
Layer} for SPARQL processing. The Application Layer is responsible for
user interaction and application state; it hosts user interfaces and
application code, and embeds Unravel sessions into web applications.
Within this layer, hooks and UI components consume Unravel queries,
while three immutable objects — \texttt{UnravelConfiguration},
\texttt{UnravelSession}, and \texttt{UnravelQuery} — encapsulate
configuration management, query construction, and query execution
respectively.

The Unravel Query Layer is implemented in Scala\footnote{https://scala-lang.org/}
and compiled to JavaScript with Scala.js\footnote{https://www.scala-js.org/}
and it provides the core query model. It is organized into several
functional subsystems: a configuration module, a core query DSL, a
high-level query API, UI integration utilities, and a variable scoping
mechanism. Together, these components provide a composable, focus-based 
abstraction over SPARQL queries that, similarly to query-based faceted search (QFS)
systems, maintains an explicit query context while supporting
incremental query construction, variable scoping, and safe composition
across application components. 
A dedicated Comunica facade\footnote{https://www.scala-js.org/libraries/facades.html}
connects this abstraction layer to the RDFJS-compatible JavaScript
ecosystem and exposes a typed interface to the underlying SPARQL
engines.

Finally, the RDF Data \& Query Layer corresponds to the existing
JavaScript/TypeScript ecosystem for RDF and SPARQL. In our deployment,
this layer is instantiated by a Comunica SPARQL engine that can query
over multiple RDF sources, including HTTP-accessible SPARQL endpoints
and local or remote RDF files. Unravel-RDF therefore focuses on bridging
application-level query construction and UI logic with these
heterogeneous RDF data sources, without reimplementing query processing.

\subsection{Integration with RDFJS and Comunica}

%% expliquer

%% Unravel DSL
%%     |
%%  typed facade
%%      |
%%  RDFJS interfaces
%%      |
%%  Comunica
%%      |
%%  SPARQL endpoints/files


\section{Programming Model}

%%%%%%%%%%%%%%%% Diagramme du modèle de focus
%%% juste après avoir introduit something, from, out, in, traverse.
\subsection{Query construction model}
\subsection{Focus-based navigation}
\subsection{Query composition and reuse}
\subsection{}

The programming model of Unravel-RDF is based on immutable query
objects and a focus variable that guides query construction, together
with a closure-based style that lets application code compose queries
incrementally.

At its core, Unravel-RDF provides a domain-specific language (DSL) for
constructing SPARQL queries. Instead of relying on raw query strings,
the DSL enables the composition of structured query fragments. Query
construction is incremental and follows a closure-based style, where
functions capture context and enable modular composition of query
pieces. All query objects are immutable, enabling predictable composition and
reuse of query fragments.

Queries are constructed relative to a focus variable, which serves as
an implicit context for navigating RDF graphs. Focus-binding primitives
(\texttt{something}, \texttt{from}) introduce or rebind this variable,
defining the scope of query construction. In contrast, graph-expansion
primitives (\texttt{out}, \texttt{in}, \texttt{traverse}) add triple
patterns relative to the current focus while updating it to the newly
reached node. This focus-based approach makes variable scoping explicit:
the current focus determines which variable is affected by each
primitive. These primitives accept both RDF IRIs and SPARQL variables as
arguments, so the same combinators can express schema-level templates
and data-level queries in a uniform way. Table~\ref{tab:query-primitives}
summarizes the core primitives and their semantics, and the example in
Figure~\ref{fig:unravel-sparql-example} shows how they can be composed
to build a query centered on a compound entity acting as the focus
variable. In this example, \texttt{session2} extends an initial query
over the focus variable \texttt{compound}, illustrating how Unravel-RDF
supports incremental refinement of SPARQL queries. In practice, these
primitives are exposed through \texttt{UnravelSession} objects, which
provide a fluent API for constructing and reusing focus-based queries.

Building on this programming model, Unravel-RDF is designed to integrate
with modern, event-driven user interfaces. On the application side, the
execution model is based on JavaScript promises combined with a
subscription mechanism, enabling asynchronous and staged processing of
queries. Results are delivered in
batches through a pagination mechanism and exposed as lazy pages that
are only fetched when requested, which avoids long-running SPARQL
requests and reduces perceived latency. These abstractions can be
consumed directly from hook-based and component-based UI frameworks
such as React, Vue, or D3-powered visualizations, bringing SPARQL-based
querying closer to modern frontend development practices while
preserving a clear separation between UI concerns, query abstraction,
and engine-level execution.

Unravel-RDF is distributed through standard web channels, including
CDN delivery and the npm ecosystem, facilitating integration into
modern JavaScript and TypeScript applications. The source code and
documentation are publicly available in the project repository
\footnote{\url{https://forge.inrae.fr/p2m2/unravel-rdf}},
enabling reuse and future extensions by other developers.

\section{Event-driven execution model}

%%%%% Diagramme d'exécution
%%% expliquer la vie d'une requete 

%%% User click
%%%      │
%%%      ▼
%%% UnravelQuery
%%%      │
%%%      ▼
%%% build SPARQL
%%%      │
%%%      ▼
%%% Comunica
%%%      │
%%%      ▼
%%% SPARQL endpoint(s)
%%%      │
%%%      ▼
%%% stream results
%%%      │
%%%      ├── progression(...)
%%%      │
%%%      ├── requestEvent(...)
%%%      │
%%%      ▼
%%% Promise.then(...)
%%%      │
%%%      ▼
%%% Update UI


%% Promise ;
%% progression ;
%% événements ;
%% pagination ;
%% cancellation éventuellement.

%% C'est ce qui distingue Unravel d'un simple query builder.

\section{Evaluation}

\subsection{FORVM DiseasesChem}

\subsection{Developer-oriented evaluation}

%% Comparer du SPARQL pure avec l equivalent Unravel

%%% Critères :

%%% lisibilité ;
%%% réutilisation ;
%%% séparation UI/query ;
%%%% évolution de la requête.


\subsection{Performance / overhead}

%% Question reviewer :
%% 
%% "Does the abstraction introduce overhead?"
%% 
%% Répondre :
%% 
%% temps génération SPARQL ;
%% temps exécution Comunica identique ;
%% coût abstraction négligeable.

\section{Discussion}

%% Tu discutes :
%% 
%% limite :
%%  - pas un remplacement SPARQL ;
%%  - pas un moteur ;
%%  - dépend de RDFJS.
%% généralisation :
%%  - autres domaines ;
%%  - autres interfaces.

\section{Conclusion}

\begin{acknowledgments}

\end{acknowledgments}

\section*{Declaration on Generative AI}

During the preparation of this work, the authors used Perplexity and
ChatGPT in order to: paraphrase and reword text, improve writing
style and clarity, and perform grammar and spelling checks. After using these
tools, the authors reviewed and edited the content as needed and take
full responsibility for the publication’s content.


%%
%% Define the bibliography file to be used
\bibliography{unravel-biblio,static-extra}

\appendix

%\section{Supplementary Material}

%\begin{figure}[ht]
%\centering
%\includegraphics[width=\linewidth]{forvm-unravel-screenshot.png}%
%\caption{Screenshot of the FORVM use case web interface built with Unravel-RDF, showing clusters of compounds, MeSH disease terms, and Planteome concepts.}
%\label{fig:forvm-screenshot}
%\end{figure}

\end{document}
%%
%% End of file
